\section[Rédaction Bourbaki n\textsuperscript{o}\,582 : les indices de Schur sur $\mathbf R$]{Rédaction Bourbaki n\textsuperscript{o}\,582 : les indices de Schur sur $\mathbf R$\protect\verifie{Redaction582}}\label{sec:bourbaki-582}

La rédaction n\textsuperscript{o}\,582 des archives Bourbaki (Archives Henri Poincaré, \href{https://archives-bourbaki.ahp-numerique.fr/items/show/1130}{item 1130}) est le rapport de Jean-Pierre Serre, daté de juin 1972 par sa notice, sur la réédition du chapitre~VIII d'\emph{Algèbre}, consacré aux modules et anneaux semi-simples. Son n\textsuperscript{o}\,3 demande que le chapitre définisse les groupes de Grothendieck de modules et explique les indices de Schur ; seule cette partie est transcrite (PDF, p.~10--12). La rédaction n\textsuperscript{o}\,596 du même chapitre (section~\ref{sec:bourbaki-596}) consacre un paragraphe à ces groupes. Sa transcription reste locale pour l'instant ; on n'en cite ici aucun passage.

\emph{L'indice de Schur.} Soient $k$ un corps, $A$ une $k$-algèbre et $S$ un $A$-module simple de rang fini sur $k$. Par le lemme de Schur, $D = \operatorname{End}_A(S)$ est un corps, en général non commutatif ; son centre $L$ est un corps commutatif qui contient $k$. On pose $m = [L : k]$, et l'on écrit $[D : L] = e^2$ : la dimension d'un corps gauche sur son centre est toujours un carré, et $e$ s'appelle l'indice de Schur. Si l'on étend les scalaires à une extension galoisienne $k'$ de $k$ assez grande, $k' \otimes_k S$ se décompose en $m$ modules simples deux à deux non isomorphes, un pour chaque plongement de $L$ dans $k'$, et chacun y apparaît $e$ fois. Le rapport énonce cette décomposition à sa page~11. En comptant les dimensions, $[D : k] = [D : L]\,[L : k] = m e^2$.

\emph{Le cas réel.} Pour $k = \mathbf R$ et $k' = \mathbf C$, le théorème de Frobenius ne laisse que trois corps de dimension finie sur $\mathbf R$ : $\mathbf R$, $\mathbf C$ et le corps $\mathbf H$ des quaternions. Chacun se rencontre déjà pour des représentations de groupes finis. Le groupe $\mathbf Z/3$, agissant sur $\mathbf R^2$ par la rotation d'un tiers de tour, a pour commutant les similitudes directes, c'est-à-dire $\mathbf C$ ; après complexification, la rotation se diagonalise avec les valeurs propres $\omega = e^{2i\pi/3}$ et $\bar\omega$, et le module se scinde en deux caractères conjugués, distincts. Le groupe des quaternions $\{\pm 1, \pm i, \pm j, \pm k\}$, agissant sur $\mathbf H$ par multiplication à gauche, a pour commutant les multiplications à droite, isomorphes à $\mathbf H$ ; après complexification, $\mathbf H \otimes_{\mathbf R} \mathbf C \simeq M_2(\mathbf C)$ est la somme de deux copies d'un même module simple de dimension~$2$.

\emph{La page.} L'exemple de la page~12 donne les trois types de modules simples réels avec leurs invariants. Pour $D = \mathbf R$, il donne $(m, e) = (1, 1)$, ce qui est juste. Pour $D = \mathbf C$ comme pour $D = \mathbf H$, le tapuscrit donne $(m, e) = (2, 2)$. Ces valeurs ne suivent pas les définitions de la page~11. La suite de la page s'accorde d'ailleurs avec les valeurs justes. Elle décrit les modules simples complexes de deuxième type, non isomorphes à leur conjugué, qui viennent de $\mathbf C$ : $T$ et $\bar T$ y apparaissent chacun une fois. Ceux de troisième type, isomorphes à leur conjugué sans être réalisables sur $\mathbf R$, viennent de $\mathbf H$ et y apparaissent deux fois ; c'est pourquoi le critère de réalisabilité qui suit leur demande une multiplicité paire.

\begin{encadre}
\begin{enonce}[Rédaction n\textsuperscript{o}\,582, n\textsuperscript{o}\,3, l'exemple de $\mathbf R$, valeurs corrigées par l'édition]
Pour $D = \mathbf C$, on a $m = 2$ et $e = 1$. Pour $D = \mathbf H$, on a $m = 1$ et $e = 2$. Aucun des deux ne donne $(m, e) = (2, 2)$.
\end{enonce}
\end{encadre}

\begin{proof}
Le corps $\mathbf C$ est commutatif : son centre est $\mathbf C$, d'où $m = [\mathbf C : \mathbf R] = 2$ et $e^2 = [\mathbf C : \mathbf C] = 1$. Un quaternion qui commute avec $i$ et avec $j$ est réel, de sorte que le centre de $\mathbf H$ est $\mathbf R$ ; d'où $m = 1$ et $e^2 = [\mathbf H : \mathbf R] = 4$. Les valeurs tapées sont exclues sans même connaître les centres : elles donneraient $[D : \mathbf R] = m e^2 = 8$, qui n'est la dimension ni de $\mathbf C$ ni de $\mathbf H$.
\end{proof}

\noindent\emph{Vérification en Lean.} Lean vérifie les dimensions qui tranchent : $[\mathbf C : \mathbf R] = 2 = 2 \cdot 1^2$ (\texttt{schur\_complex}, par \texttt{Complex.\allowbreak{}finrank\_real\_complex} de mathlib) et $[\mathbf H : \mathbf R] = 4 = 1 \cdot 2^2$ (\texttt{schur\_quaternion}, par \texttt{Quaternion.\allowbreak{}finrank\_eq\_four}). Les centres, $\mathbf C$ pour $\mathbf C$ et $\mathbf R$ pour $\mathbf H$, sont pris comme connus et ne sont pas formalisés, non plus que la relation $[D : k] = m e^2$. Les noms sont ceux de \href{\leanblobmain/Redaction582.lean}{\texttt{Redaction582.\allowbreak{}lean}}.

Ne sont pas démontrés ici : la décomposition de $k' \otimes_k S$ et l'interprétation des multiplicités par l'indice de Schur ; le critère de réalisabilité sur $\mathbf R$ d'un module semi-simple complexe ; la règle qui lit les trois types sur le caractère d'une représentation d'un groupe fini, par l'indicateur de Frobenius--Schur ; les autres demandes du n\textsuperscript{o}\,3 sur les groupes de Grothendieck, la trace et l'extension des scalaires.
